%% notes/topics/00-common/foundations.tex
%%
%% 用途:
%%   3テーマが共有している土台を記述する。
%%   個々のテーマに固有の内容はここに書かず、各テーマの
%%   definitions.tex / results.tex に置く。
%%
%% 注意:
%%   regular singular の定義は資料ごとに書き方が異なる。
%%   同値性が確認されるまでは候補を併記し、一つに断定しない。

\section{共通の土台}

\subsection{基礎設定}

\subsection{対数微分作用素}

\subsection{Frobenius 降下の塔}
%% stratified bundle / F-divided sheaf。有限レベルと無限レベルの対比。

\subsection{regular singular の定義（候補の併記）}
%% 各資料の定義を並記する。同値性は未確認事項として
%% 00-common/literature-status.tex に登録する。

\subsection{境界における不変量}
%% exponent / radius / r-スペクトラムの位置づけ。
%% 同一視の検討は 90-crosscutting/zp-invariants.tex で行う。
